% GENERATED FILE. Edit canonical lessons, then run npm run generate:books.
% canonical-source-hash: fnv1a64:0fe661a79e4b30af
% canonical-lessons: LA-C530-erus, LA-C530-verna, LA-C530-ancillula, LA-C530-servilis, LA-C530-pedisequus

\chapter{Masters and slaves}
\label{ch:la-masters-and-slaves}
\hlchaptermodality{530}

\begin{chapteropening}
\textbf{By the end of this chapter:} \emph{I can describe masters and slaves in a Roman house.}

Complete the last lesson of chapter 530: I can describe masters and slaves in a Roman house.
\end{chapteropening}

\section[erus, erī]{erus, erī — a master (of slaves)}

\label{lesson:LA-C530-erus}



\begin{quote}
Before the new one: say the Latin for a little daughter, then the Latin for brotherly.
\end{quote}

\subsection*{You'll want to know: erus, erī}
\textbf{erus, erī} — \textquotedblleft{}a master (of slaves)\textquotedblright{}.

\begin{quote}
Servus erō suō pāret.
\end{quote}

\begin{quote}
\emph{The slave obeys his master.}
\end{quote}

Masculine, second declension. The word of comedy for the master of slaves; \emph{era} is the mistress. \emph{Dominus} is the everyday word.

\begin{cousinweb}
\textbf{erus} is inherited through Proto-Italic from a Proto-Indo-European word for \textquotedblleft{}master\textquotedblright{}, and Hittite \emph{ishas} \textquotedblleft{}lord\textquotedblright{} is a cousin (EtymDB).
\end{cousinweb}

\begin{grammarlens}[title={What you've built}]
One more word for telling, explaining and describing at B1.
\end{grammarlens}

\subsection*{Your turn}
\begin{itemize}
\raggedright
  \item \emph{Say it:} \emph{erus, erī}
  \item \emph{Say it:} \emph{erus, erī}, once more
  \item \emph{Say it:} a sentence of your own with \emph{erus, erī}
  \item \emph{Recall:} say the Latin for a little daughter, then the Latin for brotherly, then say \emph{erus, erī} again
\end{itemize}

\begin{tcolorbox}[breakable,colback=teal!4,colframe=teal!35!black,title={Before you move on}]
What does erus, erī mean? (\textquotedblleft{}A master (of slaves)\textquotedblright{}.) Say it once more.
\end{tcolorbox}
\section[verna, vernae]{verna, vernae — a home-born slave}

\label{lesson:LA-C530-verna}



\begin{quote}
Before the new one: say the Latin for brotherly, then the Latin for a master.
\end{quote}

\subsection*{You'll want to know: verna, vernae}
\textbf{verna, vernae} — \textquotedblleft{}a home-born slave\textquotedblright{}.

\begin{quote}
Verna in domō dominī nātus est.
\end{quote}

\begin{quote}
\emph{A home-born slave was born in his master's house.}
\end{quote}

Masculine or feminine, first declension. Such slaves were often trusted. English \emph{vernacular} — the speech of home — comes from it.

\begin{cousinweb}
\textbf{verna} is a native Latin word whose origin is not settled; Lewis \& Short links it with a root meaning \textquotedblleft{}to dwell\textquotedblright{}. English \emph{vernacular} comes from \emph{vernāculus}, built on it.
\end{cousinweb}

\begin{grammarlens}[title={What you've built}]
One more word for telling, explaining and describing at B1.
\end{grammarlens}

\subsection*{Your turn}
\begin{itemize}
\raggedright
  \item \emph{Say it:} \emph{verna, vernae}
  \item \emph{Say it:} \emph{verna, vernae}, once more
  \item \emph{Say it:} a sentence of your own with \emph{verna, vernae}
  \item \emph{Recall:} say the Latin for brotherly, then the Latin for a master, then say \emph{verna, vernae} again
\end{itemize}

\begin{tcolorbox}[breakable,colback=teal!4,colframe=teal!35!black,title={Before you move on}]
What does verna, vernae mean? (\textquotedblleft{}A home-born slave\textquotedblright{}.) Say it once more.
\end{tcolorbox}
\section[ancillula, ancillulae]{ancillula, ancillulae — a young maidservant}

\label{lesson:LA-C530-ancillula}



\begin{quote}
Before the new one: say the Latin for a master, then the Latin for a home-born slave.
\end{quote}

\subsection*{You'll want to know: ancillula, ancillulae}
\textbf{ancillula, ancillulae} — \textquotedblleft{}a young maidservant\textquotedblright{}.

\begin{quote}
Ancillula dominae speculum dat.
\end{quote}

\begin{quote}
\emph{The little maid gives her mistress a mirror.}
\end{quote}

Feminine, a diminutive of \emph{ancilla}, already known, which is itself a diminutive.

\begin{cousinweb}
\textbf{ancillula} is a diminutive of \emph{ancilla} \textquotedblleft{}maidservant\textquotedblright{} (Lewis \& Short).
\end{cousinweb}

\begin{grammarlens}[title={What you've built}]
One more word for telling, explaining and describing at B1.
\end{grammarlens}

\subsection*{Your turn}
\begin{itemize}
\raggedright
  \item \emph{Say it:} \emph{ancillula, ancillulae}
  \item \emph{Say it:} \emph{ancillula, ancillulae}, once more
  \item \emph{Say it:} a sentence of your own with \emph{ancillula, ancillulae}
  \item \emph{Recall:} say the Latin for a master, then the Latin for a home-born slave, then say \emph{ancillula, ancillulae} again
\end{itemize}

\begin{tcolorbox}[breakable,colback=teal!4,colframe=teal!35!black,title={Before you move on}]
What does ancillula, ancillulae mean? (\textquotedblleft{}A young maidservant\textquotedblright{}.) Say it once more.
\end{tcolorbox}
\section[servīlis, servīle]{servīlis, servīle — of a slave, slavish}

\label{lesson:LA-C530-servilis}



\begin{quote}
Before the new one: say the Latin for a home-born slave, then the Latin for a young maidservant.
\end{quote}

\subsection*{You'll want to know: servīlis, servīle}
\textbf{servīlis, servīle} — \textquotedblleft{}of a slave, slavish\textquotedblright{}.

\begin{quote}
Labor servīlis dūrus est.
\end{quote}

\begin{quote}
\emph{A slave's work is hard.}
\end{quote}

Two endings, from \emph{servus}. English \emph{servile}.

\begin{cousinweb}
\textbf{servīlis} is built on \emph{servus} \textquotedblleft{}slave\textquotedblright{} (Lewis \& Short; EtymDB). English \emph{servile} comes from it.
\end{cousinweb}

\begin{grammarlens}[title={What you've built}]
One more word for telling, explaining and describing at B1.
\end{grammarlens}

\subsection*{Your turn}
\begin{itemize}
\raggedright
  \item \emph{Say it:} \emph{servīlis, servīle}
  \item \emph{Say it:} \emph{servīlis, servīle}, once more
  \item \emph{Say it:} a sentence of your own with \emph{servīlis, servīle}
  \item \emph{Recall:} say the Latin for a home-born slave, then the Latin for a young maidservant, then say \emph{servīlis, servīle} again
\end{itemize}

\begin{tcolorbox}[breakable,colback=teal!4,colframe=teal!35!black,title={Before you move on}]
What does servīlis, servīle mean? (\textquotedblleft{}Of a slave, slavish\textquotedblright{}.) Say it once more.
\end{tcolorbox}
\section[pedisequus, pedisequī]{pedisequus, pedisequī — a footman, attendant}

\label{lesson:LA-C530-pedisequus}



\begin{quote}
Before the new one: say the Latin for a young maidservant, then the Latin for of a slave.
\end{quote}

\subsection*{You'll want to know: pedisequus, pedisequī}
\textbf{pedisequus, pedisequī} — \textquotedblleft{}a footman, attendant\textquotedblright{}.

\begin{quote}
Pedisequus dominum in forum sequitur.
\end{quote}

\begin{quote}
\emph{The footman follows his master into the forum.}
\end{quote}

Masculine; \emph{pedisequa} — a female attendant. \emph{pēs} 'foot' + \emph{sequor} 'to follow'.

\begin{cousinweb}
\textbf{pedisequus} is built on \emph{pēs} \textquotedblleft{}foot\textquotedblright{} and \emph{sequor} \textquotedblleft{}to follow\textquotedblright{} (Lewis \& Short): one who follows on foot.
\end{cousinweb}

\begin{grammarlens}[title={What you've built}]
One more word for telling, explaining and describing at B1.
\end{grammarlens}

\subsection*{Your turn}
\begin{itemize}
\raggedright
  \item \emph{Say it:} \emph{pedisequus, pedisequī}
  \item \emph{Say it:} \emph{pedisequus, pedisequī}, once more
  \item \emph{Say it:} a sentence of your own with \emph{pedisequus, pedisequī}
  \item \emph{Recall:} say the Latin for a young maidservant, then the Latin for of a slave, then say \emph{pedisequus, pedisequī} again
\end{itemize}

\begin{tcolorbox}[breakable,colback=teal!4,colframe=teal!35!black,title={Before you move on}]
What does pedisequus, pedisequī mean? (\textquotedblleft{}A footman, attendant\textquotedblright{}.) Say it once more.
\end{tcolorbox}
